% Figure: parallel-plate capacitor. Uniform field between plates, charge
% on plates, fringing at the edges.
\begin{figure}[htbp]
\centering
\begin{tikzpicture}[>=Latex, scale=1.0]
% plates
\draw[engred, line width=2pt] (-2.2,1.1) -- (2.2,1.1);
\draw[engblue, line width=2pt] (-2.2,-1.1) -- (2.2,-1.1);
\node[engred, font=\small] at (2.7,1.1) {$+Q$};
\node[engblue, font=\small] at (2.7,-1.1) {$-Q$};
% uniform interior field lines
\foreach \x in {-1.8,-1.2,...,1.9}{
  \draw[enggray, ->] (\x,1.05) -- (\x,-1.05);
}
% fringing at right edge
\draw[enggray] (2.2,0.9) .. controls (2.9,0.6) and (2.9,-0.6) .. (2.2,-0.9);
\draw[enggray] (2.2,0.5) .. controls (3.3,0.2) and (3.3,-0.2) .. (2.2,-0.5);
% fringing at left edge
\draw[enggray] (-2.2,0.9) .. controls (-2.9,0.6) and (-2.9,-0.6) .. (-2.2,-0.9);
\draw[enggray] (-2.2,0.5) .. controls (-3.3,0.2) and (-3.3,-0.2) .. (-2.2,-0.5);
% spacing dimension
\draw[<->] (-2.5,1.1) -- (-2.5,-1.1) node[midway, left, font=\small] {$d$};
% area label
\node[font=\small] at (0,1.45) {plate area $A$};
\node[font=\small] at (0,0) {$E=V/d$};
\node[font=\small] at (0,-1.55) {$C=\varepsilon A/d$};
\end{tikzpicture}
\caption[Parallel-plate capacitor]{The parallel-plate capacitor. Equal
and opposite charges $\pm Q$ on plates of area $A$ at spacing $d$ set up
a uniform interior field $E = V/d$, with the field zero outside in the
ideal limit. A pillbox argument from Gauss's law gives surface charge
$\sigma = \varepsilon E$, so $Q = \varepsilon V A/d$ and the capacitance
is $C = \varepsilon A/d$. The curved fringing field at the edges is the
correction the infinite-plate formula omits; it grows in relative
importance as the gap $d$ approaches the plate dimension.}
\label{fig:vol04ch08:parallel-plate}
\end{figure}
